% ==============================================================================
% PHỤ LỤC A: DỮ LIỆU BỔ SUNG
% ==============================================================================
\chapter{Dữ liệu bổ sung}
\label{app:du-lieu}

Phụ lục chứa các bảng dài, số liệu chi tiết, mã nguồn đầy đủ hoặc biểu mẫu khảo sát.
Bảng dài dưới đây tự ngắt sang trang mới và lặp lại hàng tiêu đề.

\rowcolors{2}{white}{bkxanhnen}
\begin{longtable}{C{1.2cm} L{5cm} C{3cm} C{3cm}}
  \caption{Số liệu chi tiết theo năm (bảng dài nhiều trang)}\label{tab:bang-dai}\\
  \toprule
  \hangtieude \thead{Năm} & \thead{Chỉ tiêu} & \thead{Giá trị} & \thead{Ghi chú} \\
  \midrule
  \endfirsthead
  \toprule
  \hangtieude \thead{Năm} & \thead{Chỉ tiêu} & \thead{Giá trị} & \thead{Ghi chú} \\
  \midrule
  \endhead
  \bottomrule
  \endfoot
  2010 & Số tin tuyển dụng AI & 20.600  & -- \\
  2012 & Số tin tuyển dụng AI & 35.400  & -- \\
  2014 & Số tin tuyển dụng AI & 58.900  & -- \\
  2016 & Số tin tuyển dụng AI & 104.300 & -- \\
  2018 & Số tin tuyển dụng AI & 180.900 & tăng 8,8 lần \\
\end{longtable}
\rowcolors{1}{}{}

\section{Mã nguồn đầy đủ}

\begin{lstlisting}[language=Python, caption={Thống kê số tin theo năm}]
import pandas as pd

def thong_ke(duong_dan: str) -> pd.Series:
    """Đếm số tin tuyển dụng theo từng năm."""
    df = pd.read_csv(duong_dan, parse_dates=["ngay"])
    return df.groupby(df["ngay"].dt.year).size()
\end{lstlisting}
